% !TeX root = ../../Book4.tex
% 纲要标题：### 8.1 消解的构造
% 纲要标签：b4:sec:0701
% 纲要位置：计划纲要.md，第 3374 行。

\section{消解的构造}
\label{b4:sec:0701}

生成元给出自由满射，关系组成其核；继续为核取生成元，便得到消解。自由模可以沿满射提升映射，后面的比较定理只需要这一性质，因而可以把自由模放宽为投射对象。相反方向的构造使用内射对象的延拓性质；这两个方向都依赖范畴中相应对象是否足够多。


% 纲要内容（仅作写作计划，尚非教材正文）：
%
% * 模的自由及投射消解
% * 内射消解
% * 足够多投射或内射对象的假设
%

\subsection{模的自由及投射消解}
\label{b4:subsec:0701:01}

一个模的生成元给出自由满射，生成元之间的关系则成为这个满射的核。若核还不是容易处理的模，就继续为核选择生成元。消解把这项重复操作写成一个正合复形；它所保留的不是某一组最优生成元，而是所有关系都被逐层记录这一事实。

\begin{definition}\label{b4:def:projective-resolution}
设 \(\mathcal A\) 是交换范畴，\(P\) 为其对象。若对每个满态射 \(e:X\to Y\) 及态射 \(f:P\to Y\)，都存在态射 \(\widetilde f:P\to X\)，使 \(e\widetilde f=f\)，则称 \(P\) 为\term[tousheduixiang]{投射对象}[projective object]。给定对象 \(M\in\mathcal A\) 及增广链复形
\[
\cdots\longrightarrow P_2\xrightarrow{\mathrm{d}_2}P_1
\xrightarrow{\mathrm{d}_1}P_0\xrightarrow{\varepsilon}M\longrightarrow0
\]
若此序列正合，且每个 \(P_n\)（\(n\ge0\)）都投射，则称它为 \(M\) 的\term[toushexiaojie]{投射消解}[projective resolution]。若 \(R\) 为含幺环、\(\mathcal A=R\text{-}\mathbf{Mod}\) 为幺性左 \(R\) 模范畴，且各 \(P_n\) 自由，则称为\term[ziyouxiaojie]{自由消解}[free resolution]。
\end{definition}

正合性说明这条序列确实记录了 \(M\) 及各层关系；各项的投射性则保证从它出发的映射可以沿满态射逐层提升。后一条件将在\secref{b4:sec:0702}的比较定理中使用，不是正合性的另一个说法。

\ProseParagraphBreak
增广项 \(M\) 不属于非负次数的复形 \(P_\bullet\)：它的零次项是 \(P_0\)，从 \(P_0\) 向更低次数的微分为零。由增广诱导的商同构是
\[
H_0(P_\bullet)=P_0/\operatorname{im}\mathrm{d}_1
=P_0/\ker\varepsilon\xrightarrow{\;\overline\varepsilon\;}M.
\]
它是同构，因为 \(\varepsilon\) 满，且核正好是 \(\operatorname{im}\mathrm{d}_1\)；正次数同调则由正合性全部为零。写成上链复形时约定 \(P^{-n}=P_n\)，并在正次数补零，得到
\[
P^\bullet:\quad\cdots\longrightarrow P^{-2}\longrightarrow P^{-1}
\longrightarrow P^0\longrightarrow0.
\]
与它比较的是 \(M[0]\)，其中 \(M[0]^0=M\)，其余次数为零。增广给出复形映射 \(\varepsilon:P^\bullet\to M[0]\)：零次分量为原来的 \(\varepsilon\)，其余分量为零；\(\varepsilon\mathrm{d}^{-1}=0\) 保证复形映射条件。由刚才的同调计算，它是拟同构。

\begin{proposition}\label{b4:prop:free-resolution-exists}
设 \(R\) 是含幺环。每个幺性左 \(R\) 模都有自由消解，因而也有投射消解。每个幺性右 \(R\) 模也有相应的右模消解。
\end{proposition}
\begin{proof}
设 \(M\) 是幺性左 \(R\) 模，令 \(K_0=M\)。以 \(K_0\) 的底集合为基取自由模 \(P_0\)，把基向量 \([m]\) 送到 \(m\)，得到满射 \(p_0:P_0\to K_0\)。归纳地令 \(K_{n+1}=\ker p_n\)，取以 \(K_{n+1}\) 的底集合为基的自由模 \(P_{n+1}\) 及满射 \(p_{n+1}:P_{n+1}\to K_{n+1}\)。这里 \(K_1\) 是原生成元之间的关系模，\(K_2\) 是为 \(K_1\) 选择的生成元之间的关系模，后续各 \(K_n\) 以同样方式记录上一层的关系。

取增广 \(\varepsilon=p_0\)，对 \(n\ge1\) 定义
\[
\mathrm{d}_n:P_n\xrightarrow{p_n}K_n\hookrightarrow P_{n-1}.
\]
后一个箭头是单射，所以 \(\ker\mathrm{d}_n=\ker p_n=K_{n+1}\)；而 \(p_{n+1}\) 满射，所以 \(\operatorname{im}\mathrm{d}_{n+1}=K_{n+1}\)。因此每个正次数上像等于核，零次上也有 \(\operatorname{im}\mathrm{d}_1=K_1=\ker\varepsilon\)，增广满射，所构造的序列正合。

各项自由还保证它们投射。具体地，给定满射 \(e:X\to Y\) 及线性映射 \(f:R^{(S)}\to Y\)，按本卷的集合指标选择约定，为每个标准基向量 \(b_s\) 选择 \(x_s\in X\)，使 \(e(x_s)=f(b_s)\)。自由模的泛性质把 \(b_s\mapsto x_s\) 唯一延拓为线性映射 \(\widetilde f:R^{(S)}\to X\)；\(e\widetilde f\) 与 \(f\) 在基上相同，故 \(e\widetilde f=f\)。选择用在这一族原像上。对反环 \(R^{\mathrm{op}}\) 使用同一构造便得到右模版本。
\end{proof}

这个构造通常很大。对 \(\mathbb Z/m\mathbb Z\)，其中 \(m>1\)，较合适的消解是 \(0\to\mathbb Z\xrightarrow{\times m}\mathbb Z\to\mathbb Z/m\mathbb Z\to0\)。设 \(k\) 是域。对 \(R=k[\epsilon]/(\epsilon^2)\) 上的 \(k=R/(\epsilon)\)，则有无限消解
\[
\cdots\xrightarrow{\epsilon}R\xrightarrow{\epsilon}R
\xrightarrow{\epsilon}R\longrightarrow k\longrightarrow0.
\]
每个元素唯一写成 \(a+b\epsilon\)，其中 \(a,b\in k\)，并且 \(\epsilon(a+b\epsilon)=a\epsilon\)。因此乘 \(\epsilon\) 的核是所有 \(b\epsilon\)，像也是所有 \(a\epsilon\)，二者均为 \((\epsilon)\)；增广的核同样为 \((\epsilon)\)，所以该序列正合。这个有限生成模的消解仍有无限多个非零项，有限生成本身不使逐层取关系的过程停下来。

\subsection{内射消解}
\label{b4:subsec:0701:02}

投射性约束映射的出发对象：从它出发的映射能够沿满态射提升。内射性则约束映射的到达对象：映入它的映射能够沿单态射延拓。因而内射消解从嵌入原对象开始，逐层处理余核。

\begin{definition}\label{b4:def:injective-resolution}
设 \(\mathcal A\) 是交换范畴，\(I\) 为其对象。若对每个单态射 \(e:X\to Y\) 及态射 \(f:X\to I\)，都存在态射 \(\widetilde f:Y\to I\)，使 \(\widetilde f e=f\)，则称 \(I\) 为\term[neisheduixiang]{内射对象}[injective object]。给定对象 \(N\in\mathcal A\) 及增广上链复形
\[
0\longrightarrow N\xrightarrow{\eta}I^0\xrightarrow{\mathrm{d}^0}I^1
\xrightarrow{\mathrm{d}^1}I^2\longrightarrow\cdots
\]
若此序列正合且各 \(I^n\)（\(n\ge0\)）内射，则称它为 \(N\) 的\term[neishexiaojie]{内射消解}[injective resolution]。
\end{definition}

这里的增广项 \(N\) 不属于 \(I^\bullet\)；后者位于非负上链次数，增广 \(\eta\) 诱导 \(N\cong\ker\mathrm{d}^0=H^0(I^\bullet)\)，正次数上同调为零。构造时，先选到内射对象的单态射 \(j^0:N\hookrightarrow I^0\)，令 \(C^0=\operatorname{coker}j^0\)，再选到内射对象的单态射 \(j^1:C^0\hookrightarrow I^1\)。其余核记为 \(C^1\)，如此继续，并令
\[
\mathrm{d}^n:I^n\xrightarrow{q^n}C^n\xrightarrow{j^{n+1}}I^{n+1},
\qquad C^n=\operatorname{coker}j^n,
\]
其中 \(n\ge0\)，\(q^n\) 是余核映射，各 \(I^n\) 都内射。因为 \(j^{n+1}\) 单，\(\ker\mathrm{d}^n=\ker q^n=\operatorname{im}j^n\)；对 \(n\ge1\)，因为前一步的 \(q^{n-1}\) 满，\(\operatorname{im}\mathrm{d}^{n-1}=\operatorname{im}j^n\)。零次的核是 \(\operatorname{im}j^0\)，所以这条增广序列正合。

\ProseParagraphBreak
% 跨卷引用目标：b2:thm:enough-injectives（7.2.8）。
对任意含幺环 \(R\)，第二卷定理7.2.8已经证明每个幺性左 \(R\) 模都能嵌入内射模。这个构造先用内射交换群 \(\mathbb Q/\mathbb Z\) 分离非零元素：对 \(0\ne n\in N\)，总有加法群同态 \(\chi:N\to\mathbb Q/\mathbb Z\) 使 \(\chi(n)\ne0\)。再通过余诱导把交换群中的延拓变为 \(R\) 模中的延拓，得到内射左模
\[
J=\operatorname{Hom}_{\mathbb Z}(R,\mathbb Q/\mathbb Z),
\qquad (rf)(s)=f(sr).
\]
每个 \(\chi\) 对应左 \(R\) 线性映射 \(n\mapsto(s\mapsto\chi(sn))\)。单个映射未必分离所有元素，把全部这样的映射组成积映射，便得到嵌入
\[
N\longrightarrow\prod_{\chi\in\operatorname{Hom}_{\mathbb Z}(N,\mathbb Q/\mathbb Z)}J,
\qquad n\longmapsto\bigl(s\longmapsto\chi(sn)\bigr)_\chi.
\]
若 \(n\ne0\)，取分离它的 \(\chi\)，在 \(s=1\) 处这一分量非零，所以积映射单射。各 \(J\) 的内射性和内射模积的内射性都已在第二卷证明。本章把这个结果反复用于余核模，就得到内射消解；幺性右模的情形对反环应用同一结果。

\begin{example}\label{b4:exm:integer-injective-resolution}
作为 \(\mathbb Z\) 模，\(\mathbb Z\) 有内射消解
\[
0\longrightarrow\mathbb Z\longrightarrow\mathbb Q
\longrightarrow\mathbb Q/\mathbb Z\longrightarrow0.
\]
% 跨卷引用目标：b2:thm:divisible-injective（7.2.5）。
正合性由商群定义给出。\(\mathbb Q\) 中每个元素都能除以任意非零整数；在 \(\mathbb Q/\mathbb Z\) 中，对 \(q+\mathbb Z\) 和整数 \(m\ne0\)，有 \(m(q/m+\mathbb Z)=q+\mathbb Z\)。所以两项都可除，由第二卷定理7.2.5的可除群判据，它们都是内射模。

它与 \(\mathbb Z\) 自身的长度零投射消解适用于 Hom 的不同变量：对投射模 \(P\)，\(\operatorname{Hom}_{\mathbb Z}(P,-)\) 正合；对内射模 \(I\)，\(\operatorname{Hom}_{\mathbb Z}(-,I)\) 正合。对一般含幺环 \(R\) 及幺性左模 \(M,N\)，\secref{b4:sec:0802}将用第一变量的投射消解或第二变量的内射消解计算 \(\operatorname{Ext}_R^n(M,N)\)，并证明两种计算相同。
\end{example}

\subsection{足够多投射或内射对象的假设}
\label{b4:subsec:0701:03}

\begin{definition}\label{b4:def:enough-projectives-injectives}
设 \(\mathcal A\) 是交换范畴。若对每个对象 \(M\) 都存在投射对象 \(P\) 及满态射 \(P\twoheadrightarrow M\)，就说 \(\mathcal A\) 有\term[zugouduo-tousheduixiang]{足够多投射对象}[enough projectives]；若对每个对象 \(M\) 都存在内射对象 \(I\) 及单态射 \(M\hookrightarrow I\)，就说它有\term[zugouduo-neisheduixiang]{足够多内射对象}[enough injectives]。
\end{definition}

在前一条件下，对核反复选择投射满态射，就得到投射消解；在后一条件下，对余核反复选择内射单态射，就得到内射消解。这是\cref{b4:prop:free-resolution-exists}的范畴形式，正合性仍由每一步的核或余核保证。为每个对象选定一条消解，并不会自动得到函子：对象间的态射还需要提升到所选消解之间，提升也未必唯一。下一节的比较定理将证明提升存在且在链同伦意义下唯一；施加加性函子后再取同调，链同伦的提升给出相同的态射。这正是后续同调不变量能够成为函子的原因。

\ProseParagraphBreak
% 跨卷引用目标：b1:thm:fga-structure（9.3.1）。
交换范畴本身的公理没有包括这两项条件。例如有限交换群组成交换范畴，却没有非零投射对象。事实上，按第一卷定理9.3.1的有限生成交换群结构定理，非零有限交换群可分解为
\[
P\cong\bigoplus_i\mathbb Z/p_i^{a_i}\mathbb Z,
\qquad a_i\ge1,
\]
其中 \(p_i\) 为素数。假如 \(P\) 投射，任取一个非零直和因子 \(T=\mathbb Z/p^a\mathbb Z\)。对从 \(T\) 出发的映射，先与 \(P\to T\) 的投影复合，用 \(P\) 的投射性提升，再限制到 \(T\)，故 \(T\) 也投射。于是自然满射
\[
\mathbb Z/p^{a+1}\mathbb Z\longrightarrow\mathbb Z/p^a\mathbb Z
\]
应有截面 \(s\)。但 \(s(\bar1)\) 必须是 \(\bar1\) 的原像，可写成 \(1+p^at\pmod{p^{a+1}}\)；它不被 \(p\) 整除，所以阶为 \(p^{a+1}\)。另一方面，同态性要求
\[
p^a s(\bar1)=s(p^a\bar1)=0,
\]
所以它的阶又必须整除 \(p^a\)，矛盾。因此有限交换群范畴没有非零投射对象，不能把模范畴中的存在性结论无条件搬到它的交换子范畴。
